\section{The appointed words}

The Latin is a normalized transcription collated in the research dossier against the 1962 typical Missal and public-domain antecedents. English Scripture is the Douay–Rheims/Challoner Bible: for adapted chants it is an explicitly identified comparator, not a newly reconstructed English chant. English prayers and Preface are historical human translations for study, not an approved vernacular altar text. Abbreviated conclusions remain abbreviated. The seventeen entries include the conditional dated additions in their actual positions; the branch rules on the opening pages govern their use.

\subsection{Introit}

1962 no. 1699; printed 413 / PDF 494.

\begin{latinproper}
In voluntate tua, Domine, universa sunt posita, et non est qui possit resistere voluntati tuae: tu enim fecisti omnia, caelum et terram, et universa quae caeli ambitu continentur: Dominus universorum tu es. Ps. 118, 1. Beati immaculati in via: qui ambulant in lege Domini. V. Gloria Patri.
\end{latinproper}

Esther 13:9–11; Psalm 118:1. The following are unaltered Bible comparators, not a reconstructed English antiphon. Esther 13:9 includes the conditional Israel clause omitted by the chant; 13:11 continues beyond its last words.

\begin{englishwitness}{Historical study witness}
\textbf{13:9} And said: O Lord, Lord, almighty king, for all things are in thy power, and there is none that can resist thy will, if thou determine to save Israel.

\textbf{13:10} Thou hast made heaven and earth and all things that are under the cope of heaven.

\textbf{13:11} Thou art Lord of all, and there is none that can resist thy majesty.

\textbf{118:1} Blessed are the undefiled in the way, who walk in the law of the Lord.
\end{englishwitness}

\subsection{Sunday Collect}

1962 no. 1700; printed 414 / PDF 495.

\begin{latinproper}
Familiam tuam, quaesumus, Domine, continua pietate custodi: ut a cunctis adversitatibus, te protegente, sit libera; et in bonis actibus tuo nomini sit devota. Per Dominum nostrum Iesum Christum, Filium tuum: Qui tecum vivit et regnat in unitate Spiritus Sancti Deus: per omnia saecula saeculorum. R. Amen.
\end{latinproper}

Cummiskey 1861, XXI Sunday after Pentecost, Collect, printed 450 / PDF 459. Historical abbreviated conclusion preserved. When mission prayers are added the Sunday prayer and addition share the appropriate conclusion; the full conclusion shown here is the Sunday text standing alone.

\begin{englishwitness}{Historical study witness}
Preserve thy family, O Lord, we beseech thee, by thy constant mercy: that under thy protection, it may be freed from all adversity; and by the practice of good works, continue devoted to thy name. Thro’.
\end{englishwitness}

\subsection{Mission Collect}

1962 no. 4833; printed [86] / PDF 894.

\begin{latinproper}
Deus, qui omnes homines vis salvos fieri et ad agnitionem veritatis venire: mitte, quaesumus, operarios in messem tuam, et da eis cum omni fiducia loqui verbum tuum; ut sermo tuus currat et clarificetur, et omnes gentes cognoscant te solum Deum verum, et quem misisti Iesum Christum, Filium tuum, Dominum nostrum: Qui tecum vivit et regnat.
\end{latinproper}

Lasance–Walsh, revised 1945, printed 1453 / PDF 1474. Historical study English; mission addition under the Sunday conclusion, unless the Ordinary appoints another day.

\begin{englishwitness}{Historical study witness}
O God, Who willest that all men should be saved and come to the knowledge of the truth, send, we beseech Thee, laborers to Thy harvest, and give them to speak Thy word with all confidence, that Thy message may run and may be made plain, and that all peoples may know Thee, the only true God, and Him Whom Thou hast sent, Jesus Christ our Lord. Who with Thee.
\end{englishwitness}

\subsection{St Luke Collect}

1962 no. 3883; printed 689 / PDF 770.

\begin{latinproper}
Interveniat pro nobis, quaesumus, Domine, sanctus tuus Lucas Evangelista: qui crucis mortificationem iugiter in suo corpore, pro tui nominis honore, portavit. Per Dominum.
\end{latinproper}

Cummiskey 1861, St Luke, printed 677–678 / PDF 686–687. Historical study English; commemoration in low or conventual Masses, under a separate conclusion; absent in nonconventual sung Masses.

\begin{englishwitness}{Historical study witness}
May the holy evangelist Luke, we beseech thee, O Lord, intercede for us, who continually bore the mortification of the cross in his body for the honour of thy name. Thro’.
\end{englishwitness}

\subsection{Epistle}

1962 no. 1701; printed 414 / PDF 495.

\begin{latinproper}
Fratres: Confortamini in Domino, et in potentia virtutis eius. Induite vos armaturam Dei, ut possitis stare adversus insidias diaboli. Quoniam non est nobis colluctatio adversus carnem et sanguinem: sed adversus principes, et potestates, adversus mundi rectores tenebrarum harum, contra spiritualia nequitiae, in caelestibus. Propterea accipite armaturam Dei, ut possitis resistere in die malo, et in omnibus perfecti stare. State ergo succincti lumbos vestros in veritate, et induti loricam iustitiae, et calceati pedes in praeparatione Evangelii pacis: in omnibus sumentes scutum fidei, in quo possitis omnia tela nequissimi ignea exstinguere: et galeam salutis assumite: et gladium spiritus, quod est verbum Dei.
\end{latinproper}

Ephesians 6:10–17. DR includes “Finally, brethren” in verse 10; Latin appointed introduction is Fratres. Verse 18 is context only.

\begin{englishwitness}{Historical study witness}
\textbf{6:10} Finally, brethren, be strengthened in the Lord and in the might of his power.

\textbf{6:11} Put you on the armour of God, that you may be able to stand against the deceits of the devil.

\textbf{6:12} For our wrestling is not against flesh and blood; but against principalities and powers, against the rulers of the world of this darkness, against the spirits of wickedness in the high places.

\textbf{6:13} Therefore, take unto you the armour of God, that you may be able to resist in the evil day and to stand in all things perfect.

\textbf{6:14} Stand therefore, having your loins girt about with truth and having on the breastplate of justice:

\textbf{6:15} And your feet shod with the preparation of the gospel of peace.

\textbf{6:16} In all things taking the shield of faith, wherewith you may be able to extinguish all the fiery darts of the most wicked one.

\textbf{6:17} And take unto you the helmet of salvation and the sword of the Spirit (which is the word of God).
\end{englishwitness}

\subsection{Gradual}

1962 no. 1702; printed 414 / PDF 495.

\begin{latinproper}
Domine, refugium factus es nobis, a generatione et progenie. V. Priusquam montes fierent, aut formaretur terra et orbis: a saeculo et usque in saeculum tu es, Deus.
\end{latinproper}

Psalm 89:1–2 (Hebrew 90). The DR superscription identifies Moses but is not sung in the Gradual.

\begin{englishwitness}{Historical study witness}
\textbf{89:1} A prayer of Moses the man of God. Lord, thou hast been our refuge from generation to generation.

\textbf{89:2} Before the mountains were made, or the earth and the world was formed; from eternity and to eternity thou art God.
\end{englishwitness}

\subsection{Alleluia}

1962 no. 1703; printed 414 / PDF 495.

\begin{latinproper}
Alleluia, alleluia. V. In exitu Israel de Aegypto, domus Iacob de populo barbaro. Alleluia.
\end{latinproper}

Psalm 113:1 (Hebrew 114:1). Bible verse only; Alleluia acclamations belong to the Latin chant.

\begin{englishwitness}{Historical study witness}
\textbf{113:1} When Israel went out of Egypt, the house of Jacob from a barbarous people:
\end{englishwitness}

\subsection{Gospel}

1962 no. 1704; printed 414 / PDF 495.

\begin{latinproper}
In illo tempore: Dixit Iesus discipulis suis parabolam hanc: Assimilatum est regnum caelorum homini regi, qui voluit rationem ponere cum servis suis. Et cum coepisset rationem ponere, oblatus est ei unus, qui debebat ei decem millia talenta. Cum autem non haberet unde redderet, iussit eum dominus eius venumdari, et uxorem eius, et filios, et omnia, quae habebat, et reddi. Procidens autem servus ille, orabat eum, dicens: Patientiam habe in me, et omnia reddam tibi. Misertus autem dominus servi illius, dimisit eum, et debitum dimisit ei. Egressus autem servus ille, invenit unum de conservis suis, qui debebat ei centum denarios: et tenens suffocabat eum, dicens: Redde quod debes. Et procidens conservus eius, rogabat eum, dicens: Patientiam habe in me, et omnia reddam tibi. Ille autem noluit: sed abiit, et misit eum in carcerem, donec redderet debitum. Videntes autem conservi eius quae fiebant, contristati sunt valde: et venerunt, et narraverunt domino suo omnia, quae facta fuerant. Tunc vocavit illum dominus suus: et ait illi: Serve nequam, omne debitum dimisi tibi, quoniam rogasti me: nonne ergo oportuit et te misereri conservi tui, sicut et ego tui misertus sum? Et iratus dominus eius, tradidit eum tortoribus, quoadusque redderet universum debitum. Sic et Pater meus caelestis faciet vobis, si non remiseritis unusquisque fratri suo de cordibus vestris.
\end{latinproper}

Matthew 18:23–35. DR begins “Therefore”; the Missal substitutes its narrative introduction. The English below retains the Bible's own introduction. Credo follows by appointment.

\begin{englishwitness}{Historical study witness}
\textbf{18:23} Therefore is the kingdom of heaven likened to a king, who would take an account of his servants.

\textbf{18:24} And when he had begun to take the account, one was brought to him, that owed him ten thousand talents.

\textbf{18:25} And as he had not wherewith to pay it, his lord commanded that he should be sold, and his wife and children, and all that he had, and payment to be made.

\textbf{18:26} But that servant falling down, besought him, saying: Have patience with me, and I will pay thee all.

\textbf{18:27} And the lord of that servant being moved with pity, let him go and forgave him the debt.

\textbf{18:28} But when that servant was gone out, he found one of his fellow-servants that owed him an hundred pence: and laying hold of him, he throttled him, saying: Pay what thou owest.

\textbf{18:29} And his fellow-servant falling down, besought him, saying: Have patience with me, and I will pay thee all.

\textbf{18:30} And he would not: but went and cast him into prison, till he paid the debt.

\textbf{18:31} Now his fellow servants seeing what was done, were very much grieved, and they came, and told their lord all that was done.

\textbf{18:32} Then his lord called him: and said to him: Thou wicked servant, I forgave thee all the debt, because thou besoughtest me:

\textbf{18:33} Shouldst not thou then have had compassion also on thy fellow servant, even as I had compassion on thee?

\textbf{18:34} And his lord being angry, delivered him to the torturers until he paid all the debt.

\textbf{18:35} So also shall my heavenly Father do to you, if you forgive not every one his brother from your hearts.
\end{englishwitness}

\subsection{Offertory}

1962 no. 1705; printed 415 / PDF 496.

\begin{latinproper}
Vir erat in terra Hus, nomine Iob: simplex et rectus, ac timens Deum: quem Satan petiit, ut tentaret: et data est ei potestas a Domino in facultates et in carnem eius: perdiditque omnem substantiam ipsius et filios: carnem quoque eius gravi ulcere vulneravit.
\end{latinproper}

An adapted narrative, not consecutive “Job 1”. Complete Bible comparators below retain their own words. Job 1:13–19 supplies the losses, 2:6–7 the bodily attack. Neither these separate excerpts nor their concatenation is the appointed English chant.

\begin{englishwitness}{Historical study witness}
\textbf{1:1} There was a man in the land of Hus, whose name was Job, and that man was simple and upright, and fearing God, and avoiding evil.

\textbf{1:12} Then the Lord said to Satan: Behold, all that he hath is in thy hand: only put not forth thy hand upon his person. And Satan went forth from the presence of the Lord.

\textbf{2:6} And the Lord said to Satan: Behold, he is in thy hand, but yet save his life.

\textbf{2:7} So Satan went forth from the presence of the Lord, and struck Job with a very grievous ulcer, from the sole of the foot even to the top of his head:
\end{englishwitness}

\subsection{Sunday Secret}

1962 no. 1706; printed 415 / PDF 496.

\begin{latinproper}
Suscipe, Domine, propitius hostias: quibus et te placari voluisti, et nobis salutem potenti pietate restitui. Per Dominum nostrum.
\end{latinproper}

Cummiskey 1861, XXI Sunday after Pentecost, Secret, printed 452 / PDF 461.

\begin{englishwitness}{Historical study witness}
Mercifully receive, O Lord, these offerings, by which thou art pleased to be appeased, and restore us to salvation by thy powerful goodness. Thro’.
\end{englishwitness}

\subsection{Mission Secret}

1962 no. 4841; printed [87] / PDF 895.

\begin{latinproper}
Protector noster, aspice, Deus, et respice in faciem Christi tui, qui dedit redemptionem semetipsum pro omnibus: et fac, ut ab ortu solis usque ad occasum magnificetur nomen tuum in gentibus, ac in omni loco sacrificetur et offeratur nomini tuo oblatio munda. Per eundem Dominum.
\end{latinproper}

Lasance–Walsh, revised 1945, printed 1456 / PDF 1477. Historical study English; mission addition under the Sunday conclusion, unless the Ordinary appoints another day.

\begin{englishwitness}{Historical study witness}
Behold, O God, our protector, and look upon the face of Thy Christ, who gave Himself a redemption for all, and cause that, from the rising of the sun unto the going down thereof, Thy name be magnified among the gentiles, and that in every place a clean offering be sacrificed to Thy name. Through the same.
\end{englishwitness}

\subsection{St Luke Secret}

1962 no. 3889; printed 690 / PDF 771.

\begin{latinproper}
Donis caelestibus da nobis, quaesumus, Domine, libera tibi mente servire: ut munera quae deferimus, interveniente beato Evangelista tuo Luca, et medelam nobis operentur et gloriam. Per Dominum.
\end{latinproper}

Cummiskey 1861, St Luke, printed 677–678 / PDF 686–687. Historical study English; commemoration in low or conventual Masses, under a separate conclusion; absent in nonconventual sung Masses.

\begin{englishwitness}{Historical study witness}
Give us grace, O Almighty God, by virtue of these heavenly gifts, to serve thee with perfect liberty of mind, by the intercession of blessed Luke, thy evangelist: that the gifts which we now offer, may work in us a cure, and procure us glory. Thro’.
\end{englishwitness}

\subsection{Preface of the Most Holy Trinity}

1962 no. 1082; printed 293 / PDF 374; appointed after 1706.

\begin{latinproper}
Vere dignum et iustum est, aequum et salutare, nos tibi semper et ubique gratias agere: Domine, sancte Pater, omnipotens aeterne Deus: Qui cum unigenito Filio tuo, et Spiritu Sancto, unus es Deus, unus es Dominus: non in unius singularitate personae, sed in unius Trinitate substantiae. Quod enim de tua gloria, revelante te, credimus, hoc de Filio tuo, hoc de Spiritu Sancto, sine differentia discretionis sentimus. Ut in confessione verae sempiternaeque Deitatis, et in personis proprietas, et in essentia unitas, et in maiestate adoretur aequalitas. Quam laudant Angeli atque Archangeli, Cherubim quoque ac Seraphim: qui non cessant clamare cotidie, una voce dicentes:
\end{latinproper}

Cummiskey 1861, Trinity Preface, printed xxviii–xxix / PDF 39–40: common opening through “eternal God” followed by the Trinity continuation. The historical cue “Holy, \&c.” is retained. Dialogue and Sanctus belong to the Ordinary.

\begin{englishwitness}{Historical study witness}
It is truly meet and just, right and available to salvation, that we should always, and in all places, give thanks to thee, O holy Lord, Father Almighty, eternal God. Who together with thy only begotten Son and the Holy Ghost, art one God, and one Lord not in a singularity of one Person, but in a Trinity of one substance. For what we believe of thy glory, as thou hast revealed, the same we believe of thy Son, and of the Holy Ghost, without any difference or distinction. So that in the confession of the true and eternal Deity, we adore a distinction in the Persons, an unity in the essence, and an equality in the Majesty. Whom the angels and archangels, the cherubim also and seraphim praise: and cease not daily to cry out with one voice saying, Holy, \&c.
\end{englishwitness}

\subsection{Communion}

1962 no. 1707; printed 415 / PDF 496.

\begin{latinproper}
In salutari tuo anima mea, et in verbum tuum speravi: quando facies de persequentibus me iudicium? iniqui persecuti sunt me, adiuva me, Domine Deus meus.
\end{latinproper}

Psalm 118:81,84,86 (Hebrew 119). The Latin is selected/adapted: it omits Defecit at the opening and adds Domine Deus meus at the end. These complete DR verses are comparators, not a new translation of that adaptation. Verse 84 includes a question omitted by the chant.

\begin{englishwitness}{Historical study witness}
\textbf{118:81} My soul hath fainted after thy salvation: and in thy word I have very much hoped.

\textbf{118:84} How many are the days of thy servant: when wilt thou execute judgment on them that persecute me?

\textbf{118:86} All thy statutes are truth: they have persecuted me unjustly, do thou help me.
\end{englishwitness}

\subsection{Sunday Postcommunion}

1962 no. 1708; printed 415 / PDF 496.

\begin{latinproper}
Immortalitatis alimoniam consecuti, quaesumus, Domine: ut, quod ore percepimus, pura mente sectemur. Per Dominum nostrum.
\end{latinproper}

Cummiskey 1861, XXI Sunday after Pentecost, Postcommunion, printed 452 / PDF 461. The historical English “receive with pure souls” is retained; the commentary also considers the Latin sectemur.

\begin{englishwitness}{Historical study witness}
Having received the food of immortality, we beseech thee, O Lord, that what we have taken with our mouths we may receive with pure souls. Thro’.
\end{englishwitness}

\subsection{Mission Postcommunion}

1962 no. 4843; printed [87] / PDF 895.

\begin{latinproper}
Redemptionis nostrae munere vegetati, quaesumus, Domine: ut, hoc perpetuae salutis auxilio, fides semper vera proficiat. Per Dominum.
\end{latinproper}

Lasance–Walsh, revised 1945, printed 1457 / PDF 1478. Historical study English; mission addition under the Sunday conclusion, unless the Ordinary appoints another day.

\begin{englishwitness}{Historical study witness}
Strengthened by the gift of our redemption, we beseech Thee, O Lord, that by this aid of everlasting salvation our faith may ever increase. Through our Lord.


\end{englishwitness}

\subsection{St Luke Postcommunion}

1962 no. 3891; printed 690 / PDF 771.

\begin{latinproper}
Praesta, quaesumus, omnipotens Deus: ut, quod de sancto altari tuo accepimus, precibus beati Evangelistae tui Lucae, sanctificet animas nostras, per quod tuti esse possimus. Per Dominum.
\end{latinproper}

Cummiskey 1861, St Luke, printed 677–678 / PDF 686–687. Historical study English; commemoration in low or conventual Masses, under a separate conclusion; absent in nonconventual sung Masses.

\begin{englishwitness}{Historical study witness}
Grant, we beseech thee, O Almighty God, that what we have received from thy holy altar, may, by the prayers of blessed Luke, the evangelist, sanctify our souls, and preserve us from all danger. Thro’.


\end{englishwitness}

\subsection{What the historical English does and does not say}
Lasance's mission Collect renders \emph{clarificetur} as ``made plain'' and has no separate rendering of \emph{Filium tuum}; his Postcommunion gives ``our faith may ever increase'' for \emph{fides semper vera proficiat}. Cummiskey's Sunday Postcommunion says ``receive with pure souls'' where the Latin has \emph{pura mente sectemur}. The inherited translations are displayed unchanged; the commentary explains the Latin where the difference matters. Biblical numbering follows the Vulgate tradition: Psalms 89, 113 and 118 correspond here to Hebrew Psalms 90, 114 and 119. Full bibliographic and rights details are in the terminal references.

